% GENERATED FILE. Edit canonical lessons, then run npm run generate:books.
% canonical-source-hash: fnv1a64:e173155d9daddc80
% canonical-lessons: JA-C271-tsukamaru, JA-C271-oikakeru, JA-C271-tsuini, JA-C271-kakureru, JA-C271-tsukamaeru

\chapter{Caught at Last}
\label{ch:ja-caught-at-last}
\hlchaptermodality{271}

\begin{chapteropening}
\textbf{By the end of this chapter:} \emph{I can say how the police chased and caught someone.}

Complete the last lesson of chapter 271: I can say how the police chased and caught someone.
\end{chapteropening}

\section[\texorpdfstring{\ja{つかまる} (tsukamaru)}{つかまる (tsukamaru)}]{\ja{つかまる} (tsukamaru) — to be caught; to be arrested}

\label{lesson:JA-C271-tsukamaru}



\begin{quote}
Before the new one: say the Japanese for that (one), then the Japanese for to touch.
\end{quote}

\subsection*{You'll want to know: \ja{つかまる}}
\textbf{\ja{つかまる}} — \emph{tsukamaru} — \textquotedblleft{}to be caught; to be arrested\textquotedblright{}.

\begin{quote}
\textbf{\ja{はんにんは} \ja{えきで} \ja{つかまりました。}} — \emph{hannin wa eki de tsukamarimashita} — \textquotedblleft{}The culprit was arrested at the station.\textquotedblright{}
\end{quote}

Polite \textbf{\ja{つかまります}}, \emph{tsukamarimasu}; a godan verb. Past (polite): \textbf{\ja{はんにんが} \ja{つかまりました。}}, \emph{hannin ga tsukamarimashita}.

\begin{cousinweb}
\textbf{\ja{つかまる}} (\emph{tsukamaru}) is a native Japanese word (wago), inherited rather than borrowed from Chinese; its kanji was attached later as a kun'yomi (JMdict).
\end{cousinweb}

\begin{grammarlens}[title={What you've built}]
One more word for telling what happened.
\end{grammarlens}

\subsection*{Your turn}
\begin{itemize}
\raggedright
  \item \emph{Say it:} \emph{tsukamaru}
  \item \emph{Say it:} \emph{tsukamaru}, once more
  \item \emph{Say it:} \emph{hannin wa eki de tsukamarimashita} — The culprit was arrested at the station.
  \item \emph{Recall:} say the Japanese for that (one), then the Japanese for to touch, then say \emph{tsukamaru} again
\end{itemize}

\begin{tcolorbox}[breakable,colback=teal!4,colframe=teal!35!black,title={Before you move on}]
What does \ja{つかまる} mean? (\textquotedblleft{}To be caught; to be arrested\textquotedblright{}.) Say it once more.
\end{tcolorbox}
\section[\texorpdfstring{\ja{おいかける} (oikakeru)}{おいかける (oikakeru)}]{\ja{おいかける} (oikakeru) — to chase; to run after}

\label{lesson:JA-C271-oikakeru}



\begin{quote}
Before the new one: say the Japanese for to touch, then the Japanese for to be caught.
\end{quote}

\subsection*{You'll want to know: \ja{おいかける}}
\textbf{\ja{おいかける}} — \emph{oikakeru} — \textquotedblleft{}to chase; to run after\textquotedblright{}.

\begin{quote}
\textbf{\ja{おまわりさんが} \ja{どろぼうを} \ja{おいかけました。}} — \emph{omawarisan ga dorobō o oikakemashita} — \textquotedblleft{}A police officer ran after the thief.\textquotedblright{}
\end{quote}

Polite \textbf{\ja{おいかけます}}, \emph{oikakemasu}; an ichidan verb. Past (polite): \textbf{\ja{いぬが} \ja{ねこを} \ja{おいかけました。}}, \emph{inu ga neko o oikakemashita}.

\begin{cousinweb}
\textbf{\ja{おいかける}} (\emph{oikakeru}) is a native Japanese word (wago), inherited rather than borrowed from Chinese, built from \emph{oi}, \textquotedblleft{}chasing\textquotedblright{} (from \emph{ou}), and \emph{kakeru}, \textquotedblleft{}to set about\textquotedblright{} (JMdict, KANJIDIC2).
\end{cousinweb}

\begin{grammarlens}[title={What you've built}]
One more word for telling what happened.
\end{grammarlens}

\subsection*{Your turn}
\begin{itemize}
\raggedright
  \item \emph{Say it:} \emph{oikakeru}
  \item \emph{Say it:} \emph{oikakeru}, once more
  \item \emph{Say it:} \emph{omawarisan ga dorobō o oikakemashita} — A police officer ran after the thief.
  \item \emph{Recall:} say the Japanese for to touch, then the Japanese for to be caught, then say \emph{oikakeru} again
\end{itemize}

\begin{tcolorbox}[breakable,colback=teal!4,colframe=teal!35!black,title={Before you move on}]
What does \ja{おいかける} mean? (\textquotedblleft{}To chase; to run after\textquotedblright{}.) Say it once more.
\end{tcolorbox}
\section[\texorpdfstring{\ja{ついに} (tsuini)}{ついに (tsuini)}]{\ja{ついに} (tsuini) — finally; in the end}

\label{lesson:JA-C271-tsuini}



\begin{quote}
Before the new one: say the Japanese for to be caught, then the Japanese for to chase.
\end{quote}

\subsection*{You'll want to know: \ja{ついに}}
\textbf{\ja{ついに}} — \emph{tsuini} — \textquotedblleft{}finally; in the end\textquotedblright{}.

\begin{quote}
\textbf{\ja{ついに} \ja{はんにんが} \ja{つかまりました。}} — \emph{tsuini hannin ga tsukamarimashita} — \textquotedblleft{}At long last the culprit was caught.\textquotedblright{}
\end{quote}

\begin{cousinweb}
\textbf{\ja{ついに}} (\emph{tsuini}) is a native Japanese word (wago), inherited rather than borrowed from Chinese; its kanji was attached later as a kun'yomi (JMdict).
\end{cousinweb}

\begin{grammarlens}[title={What you've built}]
One more word for telling what happened.
\end{grammarlens}

\subsection*{Your turn}
\begin{itemize}
\raggedright
  \item \emph{Say it:} \emph{tsuini}
  \item \emph{Say it:} \emph{tsuini}, once more
  \item \emph{Say it:} \emph{tsuini hannin ga tsukamarimashita} — At long last the culprit was caught.
  \item \emph{Recall:} say the Japanese for to be caught, then the Japanese for to chase, then say \emph{tsuini} again
\end{itemize}

\begin{tcolorbox}[breakable,colback=teal!4,colframe=teal!35!black,title={Before you move on}]
What does \ja{ついに} mean? (\textquotedblleft{}Finally; in the end\textquotedblright{}.) Say it once more.
\end{tcolorbox}
\section[\texorpdfstring{\ja{かくれる} (kakureru)}{かくれる (kakureru)}]{\ja{かくれる} (kakureru) — to hide (oneself)}

\label{lesson:JA-C271-kakureru}



\begin{quote}
Before the new one: say the Japanese for to chase, then the Japanese for finally.
\end{quote}

\subsection*{You'll want to know: \ja{かくれる}}
\textbf{\ja{かくれる}} — \emph{kakureru} — \textquotedblleft{}to hide (oneself)\textquotedblright{}.

\begin{quote}
\textbf{\ja{はんにんは} \ja{やまに} \ja{かくれていました。}} — \emph{hannin wa yama ni kakurete imashita} — \textquotedblleft{}The culprit had been hiding in the hills.\textquotedblright{}
\end{quote}

Polite \textbf{\ja{かくれます}}, \emph{kakuremasu}; an ichidan verb. Past (polite): \textbf{\ja{ねこが} \ja{かくれました。}}, \emph{neko ga kakuremashita}.

\begin{cousinweb}
\textbf{\ja{かくれる}} (\emph{kakureru}) is a native Japanese word (wago), inherited rather than borrowed from Chinese; its kanji was attached later as a kun'yomi (JMdict).
\end{cousinweb}

\begin{grammarlens}[title={What you've built}]
One more word for telling what happened.
\end{grammarlens}

\subsection*{Your turn}
\begin{itemize}
\raggedright
  \item \emph{Say it:} \emph{kakureru}
  \item \emph{Say it:} \emph{kakureru}, once more
  \item \emph{Say it:} \emph{hannin wa yama ni kakurete imashita} — The culprit had been hiding in the hills.
  \item \emph{Recall:} say the Japanese for to chase, then the Japanese for finally, then say \emph{kakureru} again
\end{itemize}

\begin{tcolorbox}[breakable,colback=teal!4,colframe=teal!35!black,title={Before you move on}]
What does \ja{かくれる} mean? (\textquotedblleft{}To hide (oneself)\textquotedblright{}.) Say it once more.
\end{tcolorbox}
\section[\texorpdfstring{\ja{つかまえる} (tsukamaeru)}{つかまえる (tsukamaeru)}]{\ja{つかまえる} (tsukamaeru) — to catch; to arrest}

\label{lesson:JA-C271-tsukamaeru}



\begin{quote}
Before the new one: say the Japanese for finally, then the Japanese for to hide (oneself).
\end{quote}

\subsection*{You'll want to know: \ja{つかまえる}}
\textbf{\ja{つかまえる}} — \emph{tsukamaeru} — \textquotedblleft{}to catch; to arrest\textquotedblright{}.

\begin{quote}
\textbf{\ja{けいさつは} \ja{きのう} \ja{はんにんを} \ja{つかまえました。}} — \emph{keisatsu wa kinō hannin o tsukamaemashita} — \textquotedblleft{}The police caught the culprit yesterday.\textquotedblright{}
\end{quote}

Polite \textbf{\ja{つかまえます}}, \emph{tsukamaemasu}; an ichidan verb. Past (polite): \textbf{\ja{けいさつが} \ja{はんにんを} \ja{つかまえました。}}, \emph{keisatsu ga hannin o tsukamaemashita}.

\begin{cousinweb}
\textbf{\ja{つかまえる}} (\emph{tsukamaeru}) is a native Japanese word (wago), inherited rather than borrowed from Chinese; its kanji was attached later as a kun'yomi (JMdict). It is the transitive partner of \emph{tsukamaru}, \textquotedblleft{}to be caught\textquotedblright{}.
\end{cousinweb}

\begin{grammarlens}[title={What you've built}]
One more word for telling what happened.
\end{grammarlens}

\subsection*{Your turn}
\begin{itemize}
\raggedright
  \item \emph{Say it:} \emph{tsukamaeru}
  \item \emph{Say it:} \emph{tsukamaeru}, once more
  \item \emph{Say it:} \emph{keisatsu wa kinō hannin o tsukamaemashita} — The police caught the culprit yesterday.
  \item \emph{Recall:} say the Japanese for finally, then the Japanese for to hide (oneself), then say \emph{tsukamaeru} again
\end{itemize}

\begin{tcolorbox}[breakable,colback=teal!4,colframe=teal!35!black,title={Before you move on}]
What does \ja{つかまえる} mean? (\textquotedblleft{}To catch; to arrest\textquotedblright{}.) Say it once more.
\end{tcolorbox}
